\subsubsection{Adaptive versus Fixed Neighborhoods}
\label{subsec:adaptive_fixed_ablation}
To isolate neighborhood membership selection, the adaptive construction was replaced by fixed nearest-neighbor counts of $5$, $10$, and $20$, with $k=10$ at the base scale. Distance weighting, reliability assessment, reciprocal support, intervention thresholds, allocation, and training weights were unchanged. The adaptive reference radius was retained only for the expansion-ratio adequacy test; this comparison therefore isolates membership selection rather than removing every radius-dependent mechanism.
Both variants used identical synthetic observations, corruption realizations, and random seeds. Adaptive reruns reproduced all 2400 previously reported synthetic G-mean values exactly. Benchmark evaluation used the same partitions and corruption, with separate inner classifier selection for the fixed variant.
Across the 63 synthetic grid settings, mean G-mean was 0.7960 for adaptive and 0.7949 for fixed neighborhoods. Adaptive selection achieved higher setting-level means in 41/63 settings. These averages summarize heterogeneous conditions and do not establish uniform superiority.
Across the 17 stress settings, mean G-mean was 0.8677 for adaptive and 0.8661 for fixed neighborhoods. Adaptive selection achieved higher setting-level means in 14/17 settings. These averages summarize heterogeneous conditions and do not establish uniform superiority.
\begin{table}[t]
\caption{Adaptive minus Fixed Neighborhoods: Dataset-Level G-mean Comparisons}
\label{tab:adaptive_fixed_benchmark}
\centering
\begin{tabular}{llrrr}
\hline
Condition & Model & Mean $\Delta$ & Raw $p$ & Holm $p$\\
\hline
Supplied & LR & -0.00937 & 0.8438 & 1.0000 \\
Supplied & RF & +0.00840 & 0.0312 & 0.1875 \\
Supplied & SVM & +0.00019 & 0.8750 & 1.0000 \\
Corrupted & LR & -0.00113 & 0.5625 & 1.0000 \\
Corrupted & RF & +0.00416 & 0.2188 & 1.0000 \\
Corrupted & SVM & -0.00128 & 0.4375 & 1.0000 \\
\hline
\end{tabular}
\end{table}
Each dataset contributes its mean over 15 outer folds. Two-sided Wilcoxon tests and Holm adjustment apply to the six supplementary ablation comparisons. Individual synthetic confidence intervals are reported separately and are not simultaneous intervals.
No benchmark comparison met the 0.05 significance criterion after adjustment. Adaptation showed condition-dependent benefits rather than a consistent advantage over fixed neighborhoods.
